\documentclass[11pt]{article}
\usepackage[margin=1in]{geometry}
\usepackage{amsmath,amssymb,amsfonts,mathtools}
\usepackage{array,booktabs,enumitem,graphicx,xcolor,hyperref}
\usepackage[T1]{fontenc}
\usepackage[utf8]{inputenc}
\title{Solving 3+1 QCD on the Transverse Lattice Using 1+1 Conformal Field
  Theory}
\author{Paul A. Griffin}
\date{}
\begin{document}
\maketitle
\begin{abstract}

A new transverse lattice model of \$3+1\$ Yang-Mills theory is constructed by introducing Wess-Zumino terms into the 2-D unitary non-linear sigma model action for link fields on a 2-D lattice. The Wess-Zumino terms permit one to solve the basic non-linear sigma model dynamics of each link, for discrete values of the bare QCD coupling constant, by applying the representation theory of non-Abelian current (Kac-Moody) algebras. This construction eliminates the need to approximate the non-linear sigma model dynamics of each link with a linear sigma model theory, as in previous transverse lattice formulations. The non-perturbative behavior of the non-linear sigma model is preserved by this construction. While the new model is in principle solvable by a combination of conformal field theory, discrete light-cone, and lattice gauge theory techniques, it is more realistically suited for study with a Tamm-Dancoff truncation of excited states. In this context, it may serve as a useful framework for the study of non-perturbative phenomena in QCD via analytic techniques.

\end{abstract}

\section{Introduction}

0= 0= 0= 1= 1= 1= 2= 2= 2= = = =cmr10 scaled 1 =cmr7 scaled 1 =cmr5 scaled 1 =cmbx10 scaled 1 =cmbx7 scaled 1 =cmbx5 scaled 1 =cmmi10 scaled 1 =cmmi7 scaled 1 =cmmi5 scaled 1 =cmsy10 scaled 1 =cmsy7 scaled 1 =cmsy5 scaled 1 =cmex10 scaled 1 =cmtt10 scaled 1 =cmti10 scaled 1 =cmsl10 scaled 1

\section{Method}

0= 0= 0= 1= 1= 1= 2= 2= 2= = = = = =1pt =1pt = =24pt =0pt plus 1.2pt =12pt =4.8pt =3.6pt =14.4pt =1.2pt =0pt =13pt plus 3.6pt minus 5.8pt =13pt plus 3.6pt minus 5.8pt =-1.4pt plus 3.6pt =13pt plus 3.6pt minus 3.6pt =12pt =12pt =0.6pt =1.44pt =6pt =3.6mu =3.6mu plus 1.2mu minus 1.2mu =4mu plus 2mu minus 1mu =3.6pt plus 1.2pt minus 1.2pt =7.2pt plus 2.4pt minus 2.4pt =14.4pt plus 4.8pt minus 4.8pt

\begin{equation}
\label{eq:compact}
L(\theta) = \sum_{i=1}^{n} \ell(x_i, \theta)
\end{equation}
Equation~\ref{eq:compact} is included to keep the sample paper-like while remaining small.

\section{Conclusion}

A new transverse lattice model of math expression Yang-Mills theory is constructed by introducing Wess-Zumino terms into the 2-D unitary non-linear sigma model action for link fields on a 2-D lattice. The Wess-Zumino terms permit one to solve the basic non-linear sigma model dynamics of each link, for discrete values of the bare QCD coupling constant, by applying the representation theory of non-Abelian current (Kac-Moody) algebras. This construction eliminates the need to approximate the non-linear sigma model dynamics of each link with a linear sigma model theory, as in previous transverse lattice formulations. The non-perturbative behavior of the non-linear sigma model is preserved by this construction. While the new model is in principle solvable by a combination of conformal field theory, discrete light-cone, and lattice gauge theory techniques, it is more realistically suited for study with a Tamm-Dancoff truncation of excited states. In this context, it may serve as a useful framework for the study of non-perturbative phenomena in QCD via analytic techniques.
\end{document}
